%%% StateOfTheArt %%%
\clearpage
\section{State of the Art}

The increasing use of Deep Neural Networks in applications such as
computer vision, speech recognition, and natural language processing
has increased the demand for hardware capable of executing these
models efficiently. Although CPUs and GPUs can execute a wide range of
workloads, their general-purpose nature may not provide the required
energy efficiency for every application. This is especially relevant
in edge and embedded systems, where power consumption, area, and
memory are limited.

As a result, many domain-specific DNN accelerators have been proposed.
These accelerators target different environments, from large data
centres to small always-on sensing devices, but they usually exploit a
similar set of ideas. These include spatial parallelism, data reuse,
sparsity, and reduced numerical precision. The way in which these
techniques are combined has a direct effect on performance, power,
energy, and area \cite{sze2017efficient}.

\subsection{DNN Hardware Accelerators}

DNN accelerators can be classified according to their compute
organization, dataflow, flexibility, or implementation technology.
These categories are not exclusive, since a single accelerator may
combine ideas from several of them.

\subsubsection{Spatial and Systolic Arrays}

Many accelerators distribute multiply-accumulate operations over an
array of processing elements. Spatial architectures map different
operations and data movements directly onto this array, allowing data
to be reused close to the compute units.

Eyeriss is an early example of a spatial accelerator for convolutional
neural networks \cite{eyeriss}. It uses an array of processing elements
and a row-stationary dataflow designed to reuse weights, activations,
and partial sums. Its architecture shows the importance of the memory
hierarchy and data movement, since the cost of moving data can be as
relevant as the cost of the arithmetic operations themselves.

Systolic arrays are a more regular form of compute array. Data moves
between neighbouring processing elements in a pipelined manner, while
each element performs part of a matrix multiplication. Google's first
Tensor Processing Unit is a well-known example of this approach
\cite{tpuv1}. It was designed for neural network inference in data
centres and uses a large matrix multiplication unit supported by
software-managed on-chip memory. The TPU shows how a specialized
architecture can provide high throughput when it is designed around
the dominant operations of a workload.

\subsubsection{Precision-Scalable and Reconfigurable Accelerators}

Not all accelerators use a fixed operating point or support only one
type of workload. Some architectures introduce reconfigurability to
adapt their precision, dataflow, or power consumption.

Envision explores another part of the accelerator design space
\cite{envision}. It supports different numerical precisions and adapts
its voltage and operating frequency to trade accuracy and performance
for energy efficiency. It also exploits sparsity to avoid unnecessary
operations. This illustrates that the optimal accelerator
configuration can depend not only on the DNN model, but also on the
precision and operating conditions selected by the designer.

TinyVers targets the opposite end of the power and performance range
\cite{tinyvers}. It is an ultra-low-power system-on-chip intended for
machine learning inference at the extreme edge. It combines a RISC-V
processor, a reconfigurable machine learning accelerator, on-chip
memory, and state-retentive eMRAM. Its design shows that flexibility,
power management, and memory organization are especially important in
small always-on devices.

\subsubsection{Depth-First and Heterogeneous Accelerators}

Accelerators can also differ in the order in which layers are
executed. DepFiN uses depth-first execution and layer fusion to reduce
the external memory traffic produced by intermediate feature maps in
high-resolution convolutional networks \cite{depfin}. Instead of
completing one full layer before starting the following one, the
accelerator processes parts of several layers in sequence, allowing
intermediate data to remain closer to the compute units.

Another distinction can be made between digital accelerators and
designs based on analog or in-memory computing. Digital designs are
usually easier to control and provide predictable numerical results,
while analog approaches can offer greater compute density at the cost
of other accuracy and implementation constraints.

DIANA presents a heterogeneous approach for edge inference
\cite{diana}. It combines a conventional RISC-V host with both a
digital DNN accelerator and an analog in-memory computing core. The
digital core provides flexibility and deterministic computation,
while the analog core targets high computational density and energy
efficiency. This is an example of how different hardware
implementations can be combined to cover workloads with different
requirements.

These accelerators show that there is no single architecture that is
optimal for every DNN and every target platform. Decisions such as the
number and organization of the processing elements, the memory sizes
and bandwidths, the dataflow, the numerical precision, and the use of
sparsity create a large design space. Exploring this space manually is
difficult, which has motivated the development of automated design
space exploration tools.

\subsection{Architecture and Mapping Exploration}

Design Space Exploration tools evaluate different accelerator
configurations without requiring a complete hardware implementation
for each candidate. These tools use analytical or simulation-based
models to estimate metrics such as latency, energy, and area. This
makes it possible to discard unsuitable alternatives early in the
design process.

ZigZag is a framework for the joint exploration of DNN accelerator
architectures and mappings \cite{zigzag}. It represents the processing
array and memory hierarchy and evaluates how the operations and data
of a workload are mapped onto them. ZigZag is especially useful for
studying data reuse and the cost of data movement across the different
levels of memory.

However, this type of exploration works mainly with an abstract
description of the architecture. Parameters such as memory capacity,
bandwidth, access energy, and area describe the expected hardware, but
they do not necessarily identify the concrete hardware IP blocks that
will implement it. Therefore, an architecture that appears promising
at this level may be more difficult to implement with the components
available for a given technology.

\subsection{Automatic Accelerator Generation}

Other approaches go further and attempt to automate part of the
hardware generation process. AutoDNNChip combines a Chip Predictor
with a Chip Builder to explore accelerator configurations and produce
hardware implementations for FPGA and ASIC targets
\cite{autodnnchip}. AutoAI2C follows a similar generator-oriented
direction and focuses on automatically producing DNN accelerators for
both types of platforms \cite{autoai2c}.

These approaches reduce the manual effort required to move from a DNN
model to a hardware design. However, automatic generation also
requires stronger assumptions about the supported architecture,
target platform, and generation flow. This can make the relationship
between the exploration model and the final hardware implementation
more tightly coupled.

A different line of research uses Large Language Models to generate
or assist in the generation of HDL code. DeepRTL proposes a unified
representation for Verilog understanding and generation
\cite{deeprtl}, while RTLCoder uses an open-source dataset and a
specialized model for RTL generation \cite{liu2025rtlcoder}. These
works are relevant to hardware automation, but they address a
different problem from architectural DSE. Their main objective is the
generation of HDL descriptions rather than the systematic exploration
of accelerator architectures and hardware IP combinations.

\subsection{Relation to This Work}

The reviewed works cover several parts of the accelerator design
process. Physical accelerators such as Eyeriss, TinyVers, and DIANA
demonstrate the diversity of possible implementation decisions.
Exploration frameworks such as ZigZag allow architectures and mappings
to be evaluated before implementation. Generator-oriented approaches
such as AutoDNNChip and AutoAI2C attempt to continue the process
towards a concrete hardware design.

In the approaches reviewed here, architectural exploration and the
selection of characterized hardware IP implementations are not
treated as two separate and replaceable exploration levels. Talos
focuses on this separation. Its first level explores abstract DNN
accelerator architectures using ZigZag, while its second level maps
the selected architectures to compatible hardware IP blocks and
evaluates their PPA characteristics. This structure is intended to
keep the architectural model independent from a particular IP
library, while still considering implementation information before a
final candidate is selected.
